It is easy to see that the above discussion generalises the notion of $\lambda$-tableaux of type $\mu$ in the classical case, in the sense that a (semistandard) $\lambda$-tableau of type~$\mu$ is a~(semistandard) $\lambda$-tableau of type $(\mu|\varnothing)$. We denote by $\TT(\lambda,\mu)=\TT(\lambda,(\mu|\varnothing))$ the set of $\lambda$-tableaux of type~$\mu$ and by $\TTS(\lambda,\mu)=\TTS(\lambda,(\mu|\varnothing))$ the set of semistandard $\lambda$-tableaux of type~$\mu$. By convention, $|\TT(\varnothing,(\varnothing|\varnothing))|=1$.


\subsection{One-dimensional representations of symmetric groups}

For the remainder of this section, let $\O$ be either $\F$ or $\Z$. The signature representation of the symmetric group $\sym{n}$ over $\O$ is denoted by $\sgn \colon \sym{n} \to \{ \pm 1 \} \subseteq \O$. We shall abuse notation and also write $\sgn$ for the $\O\sym{n}$-module associated to it. In this latter context, $\sgn$ is $\O$-free of rank~$1$, with basis $\{\epsilon\}$. We denote the trivial $\O\sym{n}$-module that is $\O$-free of rank~$1$ as $\O$, with basis~$\{\bb{1}\}$. Thus, $\sigma \cdot \bb{1} = \bb{1}$ and $\sigma \cdot \epsilon = \sgn(\sigma) \epsilon$ for all $\sigma \in \sym{n}$.

When $H$ is a subgroup of $\sym{n}$, we write the respective restricted $\O H$-modules as $\O_{H}$ and $\sgn_H$. Sometimes, we shall abuse notation and write them as $\O$ and $\sgn$ when there is no confusion.

\subsection{Young permutation modules and Specht modules}

Let $\mu$ be a composition of $n$. The Young permutation module $M^{\mu}_{\O}$ is the permutation module associated to the left regular action of $\sym{n}$ on the left cosets of $\sym{\mu}$ in $\sym{n}$. In other words, $M^{\mu}_{\O} = \bigoplus_{d \sym{\mu} \in \sym{n}/\sym{\mu}} \O (d\sym{\mu})$, and $\sigma \cdot (d \sym{\mu}) = (\sigma d )\sym{\mu}$ for all $\sigma \in \sym{n}$ and $d\sym{\mu} \in \sym{n}/\sym{\mu}$.

When $\tilde{\mu}$ is a composition of $n$ obtained by rearranging some parts of $\mu$, the Young sub\-groups~$\sym{\mu}$ and~$\sym{\tilde{\mu}}$ are conjugate in $\sym{n}$, so that $M^{\mu}_{\O} \cong M^{\tilde{\mu}}_{\O}$ as $\O\sym{n}$-modules.

Let $\lambda$ be a partition. Given a $\lambda$-tableau $\t$, let $d_{\t} = \t \circ \big(\t^{\lambda}\big)^{-1}$. Then $d_{\t} \in \sym{n}$ and $d_{\t} \cdot \t^{\lambda} = d_{\t} \circ \t^{\lambda} = \t$. We define the polytabloid \begin{gather*}e_{\t}:=\sum_{\sigma \in C_{\t}} \sgn(\sigma) (\sigma \cdot (d_{\t} \sym{\lambda})) \in M^{\lambda}_{\O}.\end{gather*} The Specht module $S^{\lambda}_{\O}$ is defined to be the $\O$-submodule of $M^{\lambda}_{\O}$ spanned by $ \{ e_{\t} \colon \t \in \TT(\lambda)\}$. It is not difficult to see that $\tau \cdot e_{\t} = e_{\tau \cdot \t}$ for $\tau \in \sym{n}$ and $\t \in \TT(\lambda)$, so that $S^{\lambda}_{\O}$ is an $\O \sym{n}$-submodule of $M^{\lambda}_{\O}$.

Let $\O=\F$ and $M^*$ denote the contragradient dual of an $\F \sym{n}$-module $M$. The following isomorphism is well known (see \cite[Theorem~8.15]{GJ}):
\begin{gather*}S^\lambda_\F\otimes \sgn\cong \big(S^{\lambda'}_\F\big)^*.\end{gather*}
Furthermore, $\F\sym{n}$ is semisimple as an algebra if and only if either $\Char(\F)=0$ or $\Char(\F)>n$, in which case, the Specht modules $S^{\lambda}_\F$, as $\lambda$ runs over all the partitions of $n$, give a complete list of pairwise non-isomorphic irreducible $\F\sym{n}$-modules.

\subsection{Signed permutation modules} \label{S:signedpermutation}

Let $(\alpha|\beta)$ be a bicomposition of $n$. We define the signed permutation module $M_{\O}(\alpha|\beta)$ to be
\begin{gather*}
M_\O(\alpha|\beta) = \Ind_{\sym{\alpha} \times \sym{\beta}}^{\sym{n}} (\O \boxtimes \sgn),
\end{gather*}
where $\sym{\alpha} \times \sym{\beta}$ is identified with the Young subgroup $\sym{\alpha|\beta}$ of $\sym{n}$. Thus, when $\Gamma$ is a left transversal of $\sym{\alpha|\beta}$ in $\sym{n}$, $M_\O(\alpha|\beta)$ has a basis $\{ d \otimes \bb{1} \otimes \epsilon \colon d \in \Gamma \}$, and $\sigma \cdot (d \otimes \bb{1} \otimes \epsilon) = \sgn(\xi_{\beta}) (d' \otimes \bb{1} \otimes \epsilon)$ if $\sigma d = d' \xi_{\alpha} \tlxi{\beta}{|\alpha|}$ where $d' \in \Gamma$, $(\xi_{\alpha}, \xi_{\beta}) \in \sym{\alpha} \times \sym{\beta}$.

Observe that if $\tilde{\alpha}$ and $\tilde{\beta}$ are compositions obtained by rearranging some parts of $\alpha$ and $\beta$ respectively, then $M_\O(\tilde{\alpha}|\tilde{\beta}) \cong M_\O(\alpha|\beta)$. Also,
\begin{gather*}
M_\O(\alpha|\beta)\otimes \sgn \cong \Ind_{\sym{\alpha}\times \sym{\beta}}^{\sym{n}}((\O\otimes\sgn)\boxtimes (\sgn\otimes \sgn))\cong M_\O(\beta|\alpha).
\end{gather*}
